\documentclass[10pt,a4paper]{article}
\usepackage[utf8]{inputenc}
\usepackage[french]{babel}
\usepackage{amsmath,amsfonts,amssymb}
\usepackage{geometry}
\usepackage{multicol}
\usepackage{xcolor}
\usepackage{tikz}
\usepackage{fancyhdr}

\geometry{margin=1.5cm}
\pagestyle{fancy}
\fancyhf{}
\fancyhead[C]{Fiche de Révision - Probabilités Numériques}
\fancyfoot[C]{\thepage}

\setlength{\columnsep}{1cm}
\setlength{\parindent}{0pt}

\title{\vspace{-2cm}\textbf{FICHE DE RÉVISION}\\
\large Probabilités Numériques - Inégalités, Notions et Astuces}
\author{}
\date{}

\begin{document}

\maketitle
\begin{multicols}{2}

\section*{INÉGALITÉS FONDAMENTALES}

\subsection*{Inégalité de Markov}
Pour $X \geq 0$ et $a > 0$ :
\[
\mathbb{P}(X \geq a) \leq \frac{\mathbb{E}[X]}{a}
\]

\subsection*{Inégalité de Tchebychev}
Pour toute v.a. $X$ avec $\mathbb{E}[X] = \mu$, $\text{Var}(X) = \sigma^2$ :
\[
\mathbb{P}(|X - \mu| \geq \varepsilon) \leq \frac{\sigma^2}{\varepsilon^2}
\]

\subsection*{Inégalité de Bienaymé-Tchebychev}
Version généralisée :
\[
\mathbb{P}(|X - \mu| \geq k\sigma) \leq \frac{1}{k^2}
\]

\subsection*{Inégalité de Jensen}
Si $\varphi$ est convexe et $X$ intégrable :
\[
\varphi(\mathbb{E}[X]) \leq \mathbb{E}[\varphi(X)]
\]

\subsection*{Inégalité de Hölder}
Pour $p, q > 1$ avec $\frac{1}{p} + \frac{1}{q} = 1$ :
\[
\mathbb{E}[|XY|] \leq (\mathbb{E}[|X|^p])^{1/p} (\mathbb{E}[|Y|^q])^{1/q}
\]

\subsection*{Inégalité de Minkowski}
Pour $p \geq 1$ :
\[
(\mathbb{E}[|X+Y|^p])^{1/p} \leq (\mathbb{E}[|X|^p])^{1/p} + (\mathbb{E}[|Y|^p])^{1/p}
\]

\subsection*{Inégalité de Cauchy-Schwarz}
Cas particulier de Hölder avec $p = q = 2$ :
\[
|\mathbb{E}[XY]| \leq \sqrt{\mathbb{E}[X^2] \mathbb{E}[Y^2]}
\]

\subsection*{Inégalité de Chernoff}
Pour $X_1, \ldots, X_n$ i.i.d. avec $\mathbb{E}[X_i] = 0$ :
\[
\mathbb{P}\left(\sum_{i=1}^n X_i \geq t\right) \leq \inf_{\lambda > 0} e^{-\lambda t} \prod_{i=1}^n \mathbb{E}[e^{\lambda X_i}]
\]

\subsection*{Inégalité de Bernstein}
Pour $X_1, \ldots, X_n$ i.i.d. avec $|X_i| \leq M$ p.s. :
\[
\mathbb{P}\left(\left|\sum_{i=1}^n X_i\right| \geq t\right) \leq 2 \exp\left(-\frac{t^2/2}{\sum_{i=1}^n \text{Var}(X_i) + Mt/3}\right)
\]

\subsection*{Inégalité de Hoeffding}
Pour $X_1, \ldots, X_n$ indépendantes avec $a_i \leq X_i \leq b_i$ :
\[
\mathbb{P}\left(\sum_{i=1}^n (X_i - \mathbb{E}[X_i]) \geq t\right) \leq \exp\left(-\frac{2t^2}{\sum_{i=1}^n (b_i - a_i)^2}\right)
\]

\section*{THÉORÈMES DE CONVERGENCE}

\subsection*{Loi faible des grands nombres}
Si $(X_n)$ i.i.d. avec $\mathbb{E}[X_1] = \mu$, $\text{Var}(X_1) < \infty$ :
\[
\frac{1}{n}\sum_{i=1}^n X_i \xrightarrow[n \to \infty]{\mathbb{P}} \mu
\]

\subsection*{Loi forte des grands nombres}
Si $(X_n)$ i.i.d. avec $\mathbb{E}[|X_1|] < \infty$ :
\[
\frac{1}{n}\sum_{i=1}^n X_i \xrightarrow[n \to \infty]{\text{p.s.}} \mu
\]

\subsection*{Théorème central limite}
Si $(X_n)$ i.i.d. avec $\mathbb{E}[X_1] = \mu$, $\text{Var}(X_1) = \sigma^2 \in (0, \infty)$ :
\[
\frac{\sqrt{n}(\bar{X}_n - \mu)}{\sigma} \xrightarrow[n \to \infty]{\mathcal{L}} \mathcal{N}(0,1)
\]

\subsection*{Théorème de Slutsky}
Si $X_n \xrightarrow{\mathcal{L}} X$ et $Y_n \xrightarrow{\mathbb{P}} c$ :
\[
X_n + Y_n \xrightarrow{\mathcal{L}} X + c, \quad X_n Y_n \xrightarrow{\mathcal{L}} cX
\]

\subsection*{Méthode delta}
Si $\sqrt{n}(T_n - \theta) \xrightarrow{\mathcal{L}} \mathcal{N}(0, \sigma^2)$ et $g$ dérivable :
\[
\sqrt{n}(g(T_n) - g(\theta)) \xrightarrow{\mathcal{L}} \mathcal{N}(0, (g'(\theta))^2 \sigma^2)
\]

\section*{MÉTHODES MONTE CARLO}

\subsection*{Estimateur classique}
\[
\hat{I}_n = \frac{1}{n} \sum_{i=1}^n f(U_i), \quad U_i \sim \mathcal{U}(0,1)
\]

\subsection*{Variance}
\[
\text{Var}(\hat{I}_n) = \frac{1}{n} \left(\int_0^1 f^2(x) dx - I^2\right) = \frac{\sigma^2}{n}
\]

\subsection*{Intervalle de confiance}
\[
\hat{I}_n \pm z_{1-\alpha/2} \frac{\hat{\sigma}_n}{\sqrt{n}}
\]

\subsection*{Échantillonnage préférentiel}
\[
\hat{I}_n^{IS} = \frac{1}{n} \sum_{i=1}^n \frac{f(X_i)}{g(X_i)}, \quad X_i \sim g
\]

\subsection*{Densité optimale}
\[
g^*(x) = \frac{|f(x)|}{\int |f(t)| dt}
\]

\section*{RÉDUCTION DE VARIANCE}

\subsection*{Variables antithétiques}
\[
\hat{I}_n^{AV} = \frac{1}{n} \sum_{i=1}^{n/2} \frac{f(U_i) + f(1-U_i)}{2}
\]
Efficace si $f$ monotone.

\subsection*{Variables de contrôle}
\[
\hat{I}_n^{CV} = \frac{1}{n} \sum_{i=1}^n f(X_i) - c \left(\frac{1}{n} \sum_{i=1}^n Y_i - \mu_Y\right)
\]
Coefficient optimal : $c^* = \frac{\text{Cov}(f(X), Y)}{\sigma_Y^2}$

\subsection*{Échantillonnage stratifié}
\[
\hat{I}_n^{ST} = \sum_{j=1}^m p_j \cdot \frac{1}{n_j} \sum_{i=1}^{n_j} f(X_{j,i})
\]

\subsection*{Allocation optimale}
\[
n_j^* = n \frac{p_j \sigma_j}{\sum_{k=1}^m p_k \sigma_k}
\]

\section*{GÉNÉRATION DE NOMBRES ALÉATOIRES}

\subsection*{Méthode de la transformée inverse}
Si $U \sim \mathcal{U}(0,1)$ et $F$ continue strictement croissante :
\[
X = F^{-1}(U) \sim F
\]

\subsection*{Exemples}
\begin{itemize}
    \item Exponentielle : $X = -\frac{1}{\lambda} \ln(1-U)$
    \item Cauchy : $X = \tan(\pi(U - 1/2))$
    \item Pareto : $X = x_m (1-U)^{-1/\alpha}$
\end{itemize}

\subsection*{Méthode de rejet}
Générer $X \sim g$, $U \sim \mathcal{U}(0,1)$.
Accepter si $U \leq \frac{f(X)}{M g(X)}$ où $f \leq M g$.

Probabilité d'acceptation : $\frac{1}{M}$.

\subsection*{Box-Muller}
Pour $U_1, U_2 \sim \mathcal{U}(0,1)$ indépendantes :
\begin{align*}
    Z_1 &= \sqrt{-2 \ln U_1} \cos(2\pi U_2) \\
    Z_2 &= \sqrt{-2 \ln U_1} \sin(2\pi U_2)
\end{align*}
Alors $Z_1, Z_2 \sim \mathcal{N}(0,1)$ indépendantes.

\section*{ESTIMATION}

\subsection*{Maximum de vraisemblance}
\[
\hat{\theta}_n = \arg\max_{\theta} \sum_{i=1}^n \ln f_\theta(X_i)
\]

\subsection*{Information de Fisher}
\[
I(\theta) = \mathbb{E}_\theta\left[\left(\frac{\partial \ln f_\theta(X)}{\partial \theta}\right)^2\right]
\]

\subsection*{Borne de Cramér-Rao}
\[
\text{Var}(\hat{\theta}_n) \geq \frac{1}{n I(\theta)}
\]

\subsection*{Normalité asymptotique MLE}
\[
\sqrt{n}(\hat{\theta}_n - \theta_0) \xrightarrow{\mathcal{L}} \mathcal{N}(0, I(\theta_0)^{-1})
\]

\section*{TESTS D'HYPOTHÈSES}

\subsection*{Test du rapport de vraisemblance}
\[
\Lambda_n = 2(\ell_n(\hat{\theta}_n) - \ell_n(\theta_0)) \xrightarrow{\mathcal{L}} \chi^2(p)
\]

\subsection*{Test de Wald}
\[
W_n = n(\hat{\theta}_n - \theta_0)^T I(\hat{\theta}_n) (\hat{\theta}_n - \theta_0) \xrightarrow{\mathcal{L}} \chi^2(p)
\]

\subsection*{Test du score}
\[
S_n = \frac{1}{n} \left(\frac{\partial \ell_n}{\partial \theta}(\theta_0)\right)^T I(\theta_0)^{-1} \left(\frac{\partial \ell_n}{\partial \theta}(\theta_0)\right) \xrightarrow{\mathcal{L}} \chi^2(p)
\]

\section*{BOOTSTRAP}

\subsection*{Variance bootstrap}
\[
\widehat{\text{Var}}(\hat{\theta}_n) = \frac{1}{B-1} \sum_{b=1}^B (\hat{\theta}_{n,b}^* - \bar{\theta}_n^*)^2
\]

\subsection*{Intervalle percentile}
\[
[\hat{\theta}_{n,(\alpha/2)}^*, \hat{\theta}_{n,(1-\alpha/2)}^*]
\]

\subsection*{Bootstrap-t}
\[
T_n^* = \frac{\hat{\theta}_n^* - \hat{\theta}_n}{\hat{\sigma}_n^*}
\]

\section*{MCMC}

\subsection*{Metropolis-Hastings}
Acceptation avec probabilité :
\[
\alpha = \min\left(1, \frac{\pi(Y) q(X|Y)}{\pi(X) q(Y|X)}\right)
\]

\subsection*{Gibbs}
Mise à jour coordonnée par coordonnée :
\[
X_i^{(n)} \sim \pi(\cdot | X_1^{(n)}, \ldots, X_{i-1}^{(n)}, X_{i+1}^{(n-1)}, \ldots, X_d^{(n-1)})
\]

\section*{ASTUCES ET TRICKS}

\subsection*{Calcul de variance}
\[
\text{Var}(X) = \mathbb{E}[X^2] - (\mathbb{E}[X])^2
\]

\subsection*{Variance d'une somme}
\[
\text{Var}\left(\sum_{i=1}^n X_i\right) = \sum_{i=1}^n \text{Var}(X_i) + 2 \sum_{i<j} \text{Cov}(X_i, X_j)
\]

Si indépendantes : $\text{Var}\left(\sum_{i=1}^n X_i\right) = \sum_{i=1}^n \text{Var}(X_i)$

\subsection*{Covariance}
\[
\text{Cov}(X, Y) = \mathbb{E}[XY] - \mathbb{E}[X]\mathbb{E}[Y]
\]

\subsection*{Loi conditionnelle}
\[
f_{X|Y}(x|y) = \frac{f_{X,Y}(x,y)}{f_Y(y)}
\]

\subsection*{Espérance conditionnelle}
\[
\mathbb{E}[X] = \mathbb{E}[\mathbb{E}[X|Y]]
\]

\subsection*{Variance totale}
\[
\text{Var}(X) = \mathbb{E}[\text{Var}(X|Y)] + \text{Var}(\mathbb{E}[X|Y])
\]

\subsection*{Moment générateur}
\[
M_X(t) = \mathbb{E}[e^{tX}]
\]
Si $M_X(t) = M_Y(t)$ pour tout $t$, alors $X \stackrel{d}{=} Y$.

\subsection*{Fonction caractéristique}
\[
\phi_X(t) = \mathbb{E}[e^{itX}]
\]
Détermine uniquement la loi.

\subsection*{Méthode des moments}
Résoudre $\mathbb{E}_\theta[X^k] = \frac{1}{n}\sum_{i=1}^n X_i^k$ pour $k = 1, 2, \ldots$

\subsection*{Estimateur sans biais}
\[
\mathbb{E}[\hat{\theta}_n] = \theta
\]

\subsection*{Biais}
\[
\text{Biais}(\hat{\theta}_n) = \mathbb{E}[\hat{\theta}_n] - \theta
\]

\subsection*{Erreur quadratique moyenne}
\[
\text{EQM}(\hat{\theta}_n) = \text{Var}(\hat{\theta}_n) + (\text{Biais}(\hat{\theta}_n))^2
\]

\subsection*{Efficacité relative}
\[
\text{ER}(\hat{\theta}_1, \hat{\theta}_2) = \frac{\text{Var}(\hat{\theta}_2)}{\text{Var}(\hat{\theta}_1)}
\]

\subsection*{Convergence en probabilité}
\[
X_n \xrightarrow{\mathbb{P}} X \Leftrightarrow \forall \varepsilon > 0, \lim_{n \to \infty} \mathbb{P}(|X_n - X| > \varepsilon) = 0
\]

\subsection*{Convergence presque sûre}
\[
X_n \xrightarrow{\text{p.s.}} X \Leftrightarrow \mathbb{P}(\lim_{n \to \infty} X_n = X) = 1
\]

\subsection*{Convergence en loi}
\[
X_n \xrightarrow{\mathcal{L}} X \Leftrightarrow \forall x \text{ point de continuité de } F_X, \lim_{n \to \infty} F_{X_n}(x) = F_X(x)
\]

\subsection*{Théorème de continuité de Lévy}
\[
X_n \xrightarrow{\mathcal{L}} X \Leftrightarrow \phi_{X_n}(t) \to \phi_X(t) \text{ pour tout } t
\]

\subsection*{Loi des grands nombres pour fonctions}
Si $X_n \xrightarrow{\text{p.s.}} X$ et $g$ continue :
\[
g(X_n) \xrightarrow{\text{p.s.}} g(X)
\]

\subsection*{Méthode de la fonction indicatrice}
\[
\mathbb{E}[\mathbf{1}_A] = \mathbb{P}(A)
\]

\subsection*{Changement de variable}
Si $Y = g(X)$ avec $g$ bijective :
\[
f_Y(y) = f_X(g^{-1}(y)) \left|\frac{d}{dy} g^{-1}(y)\right|
\]

\subsection*{Formule de convolution}
Pour $Z = X + Y$ indépendantes :
\[
f_Z(z) = \int_{-\infty}^{\infty} f_X(x) f_Y(z-x) dx
\]

\subsection*{Méthode de la fonction génératrice}
Pour variables discrètes :
\[
G_X(s) = \mathbb{E}[s^X] = \sum_{k=0}^{\infty} s^k \mathbb{P}(X = k)
\]

\subsection*{Trick : Symétrie}
Pour $U \sim \mathcal{U}(0,1)$, on a $1-U \sim \mathcal{U}(0,1)$ aussi.

\subsection*{Trick : Normalisation}
Pour estimer $\int f(x) dx$, on peut écrire :
\[
\int f(x) dx = \int \frac{f(x)}{g(x)} g(x) dx = \mathbb{E}_g\left[\frac{f(X)}{g(X)}\right]
\]

\subsection*{Trick : Importance sampling}
Choisir $g$ proche de $f$ pour réduire la variance.

\subsection*{Trick : Stratification}
Diviser l'espace en régions de variance similaire.

\subsection*{Trick : Variables antithétiques}
Utiliser $U$ et $1-U$ pour créer corrélation négative.

\subsection*{Trick : Variables de contrôle}
Soustraire une variable de moyenne connue corrélée.

\subsection*{Trick : Rao-Blackwell}
Si $T$ est suffisant et $\hat{\theta}$ estimateur :
\[
\mathbb{E}[\hat{\theta}|T] \text{ est meilleur que } \hat{\theta}
\]

\subsection*{Trick : Méthode delta multivariée}
Si $\sqrt{n}(\mathbf{T}_n - \boldsymbol{\theta}) \xrightarrow{\mathcal{L}} \mathcal{N}(\mathbf{0}, \boldsymbol{\Sigma})$ :
\[
\sqrt{n}(g(\mathbf{T}_n) - g(\boldsymbol{\theta})) \xrightarrow{\mathcal{L}} \mathcal{N}(0, \nabla g(\boldsymbol{\theta})^T \boldsymbol{\Sigma} \nabla g(\boldsymbol{\theta}))
\]

\subsection*{Trick : Loi conditionnelle normale}
Si $(X, Y) \sim \mathcal{N}(\boldsymbol{\mu}, \boldsymbol{\Sigma})$ :
\[
X|Y = y \sim \mathcal{N}(\mu_X + \Sigma_{XY} \Sigma_{YY}^{-1}(y - \mu_Y), \Sigma_{XX} - \Sigma_{XY} \Sigma_{YY}^{-1} \Sigma_{YX})
\]

\subsection*{Trick : Propriété de Markov}
Pour chaîne de Markov :
\[
\mathbb{P}(X_{n+1} = x_{n+1} | X_0, \ldots, X_n) = \mathbb{P}(X_{n+1} = x_{n+1} | X_n)
\]

\subsection*{Trick : Équation de Chapman-Kolmogorov}
\[
P^{(n+m)}_{ij} = \sum_k P^{(n)}_{ik} P^{(m)}_{kj}
\]

\subsection*{Trick : Condition de réversibilité}
Pour distribution stationnaire $\pi$ :
\[
\pi_i P_{ij} = \pi_j P_{ji} \quad \text{(balance détaillée)}
\]

\subsection*{Trick : Burn-in en MCMC}
Éliminer les premières itérations pour laisser la chaîne converger.

\subsection*{Trick : Thinning en MCMC}
Ne garder qu'une itération sur $k$ pour réduire l'autocorrélation.

\subsection*{Trick : Diagnostic de convergence}
Utiliser plusieurs chaînes avec points de départ différents.

\subsection*{Trick : Facteur de réduction potentiel}
\[
\hat{R} = \sqrt{\frac{\text{Var entre chaînes} + \text{Var dans chaînes}}{\text{Var dans chaînes}}}
\]
On veut $\hat{R} < 1.1$.

\section*{LOIS USUELLES}

\subsection*{Normale}
$X \sim \mathcal{N}(\mu, \sigma^2)$ :
\[
f(x) = \frac{1}{\sigma\sqrt{2\pi}} e^{-\frac{(x-\mu)^2}{2\sigma^2}}
\]
$\mathbb{E}[X] = \mu$, $\text{Var}(X) = \sigma^2$

\subsection*{Exponentielle}
$X \sim \text{Exp}(\lambda)$ :
\[
f(x) = \lambda e^{-\lambda x} \mathbf{1}_{x \geq 0}
\]
$\mathbb{E}[X] = \frac{1}{\lambda}$, $\text{Var}(X) = \frac{1}{\lambda^2}$

\subsection*{Uniforme}
$X \sim \mathcal{U}(a, b)$ :
\[
f(x) = \frac{1}{b-a} \mathbf{1}_{[a,b]}(x)
\]
$\mathbb{E}[X] = \frac{a+b}{2}$, $\text{Var}(X) = \frac{(b-a)^2}{12}$

\subsection*{Gamma}
$X \sim \Gamma(\alpha, \beta)$ :
\[
f(x) = \frac{\beta^\alpha}{\Gamma(\alpha)} x^{\alpha-1} e^{-\beta x} \mathbf{1}_{x > 0}
\]
$\mathbb{E}[X] = \frac{\alpha}{\beta}$, $\text{Var}(X) = \frac{\alpha}{\beta^2}$

\subsection*{Beta}
$X \sim \text{Beta}(\alpha, \beta)$ :
\[
f(x) = \frac{\Gamma(\alpha+\beta)}{\Gamma(\alpha)\Gamma(\beta)} x^{\alpha-1} (1-x)^{\beta-1} \mathbf{1}_{[0,1]}(x)
\]
$\mathbb{E}[X] = \frac{\alpha}{\alpha+\beta}$, $\text{Var}(X) = \frac{\alpha\beta}{(\alpha+\beta)^2(\alpha+\beta+1)}$

\subsection*{Chi-deux}
$X \sim \chi^2(k)$ :
\[
X = \sum_{i=1}^k Z_i^2, \quad Z_i \sim \mathcal{N}(0,1) \text{ i.i.d.}
\]
$\mathbb{E}[X] = k$, $\text{Var}(X) = 2k$

\subsection*{Student}
$T \sim t(k)$ :
\[
T = \frac{Z}{\sqrt{Y/k}}, \quad Z \sim \mathcal{N}(0,1), Y \sim \chi^2(k)
\]

\subsection*{Fisher}
$F \sim F(d_1, d_2)$ :
\[
F = \frac{X_1/d_1}{X_2/d_2}, \quad X_1 \sim \chi^2(d_1), X_2 \sim \chi^2(d_2)
\]

\end{multicols}

\end{document}
